Für die Kompression gilt $\Delta V = -\chi\,V\,\Delta p$ (FoTaBe S.~164). Gleicher Bruchteil $\Delta V/V$ bedeutet
\begin{align}
 \chi_1\,\Delta p_1 &= \chi_2\,\Delta p_2 \quad\Rightarrow\quad \frac{\chi_2}{\chi_1} = \frac{\Delta p_1}{\Delta p_2}
\end{align}
Die Kompressibilität ist umgekehrt proportional zum nötigen Überdruck. Braucht die zweite Flüssigkeit einen um \dprO\ grösseren Überdruck, $\Delta p_2 = \left(1+\dprO\right)\Delta p_1$, so ist ihre Kompressibilität um
\begin{align}
 \frac{\chi_1-\chi_2}{\chi_1} &= \rkbF = 1-\frac{1}{1+\dprO} = \rkb \approx \result{\rkbP{0}{2}}
\end{align}
kleiner. Umgekehrt ist die Kompressibilität der ersten Flüssigkeit um \dprO\ grösser als die der zweiten.
